\section{Cote 134-8, Lemma 2: a complete counit criterion}
\label{sec:SM-cote1348-counit}

The manuscript at handwritten pages 52--53 considers an adjunction from
an abelian category to a category of complexes, the inverse image of
quasi-isomorphisms, and tests on the adjunction counit.  Its conditions
labelled $(ii')$ and $(ii'')$ have revisions that should not be
silently identified with a modern assertion.  The following gives a
precise version of the clear counit/projective-object criterion.  All
complexes here are homological and concentrated in degrees at least
zero.  The exactness hypothesis on the right adjoint is explicit.

\paragraph{Finding F-SM-17 (localization equivalence and projective counit tests).}
Let $\mathcal A$ and $\mathcal M$ be abelian categories, with
$\mathcal A$ having enough projectives.  Put
$\mathcal B=\operatorname{Ch}_{\geq0}(\mathcal A)$ and let
\[
 \Lambda:\mathcal M\rightleftarrows\mathcal B:G
 \qquad(\Lambda\dashv G)
\]
be an additive adjunction such that both functors are exact.  Write
$\eta$ and $\varepsilon$ for its unit and counit, let
$\mathcal Q$ be the quasi-isomorphisms in $\mathcal B$, and put
$W=\Lambda^{-1}(\mathcal Q)$.  Fix universes in which the following
localizations exist.  Then these conditions are equivalent:
\begin{enumerate}
\item For every $C\in\mathcal B$, the counit
      $\varepsilon_C:\Lambda G(C)\to C$ is a quasi-isomorphism.
\item The same assertion holds for every projective object $P$ of
      the abelian category $\mathcal B$.
\item $G$ sends $\mathcal Q$ into $W$, and the induced functor
      \[
       \overline\Lambda:\mathcal M[W^{-1}]
          \longrightarrow \mathcal B[\mathcal Q^{-1}]
          =D_{\geq0}(\mathcal A)
      \]
      is an equivalence.
\end{enumerate}
Whenever these hold, $G^{-1}(W)=\mathcal Q$, the induced functors
$\overline\Lambda\dashv\overline G$ are mutually inverse equivalences,
and every $\eta_X$ belongs to $W$.

The equivalence of the first and third conditions does not require
abelian categories or exactness; it uses the displayed adjunction and
the saturation of quasi-isomorphisms.  For the implication from the
second condition to the first, exactness of the composite
$H=\Lambda G$ suffices in place of separate exactness of both functors.
No preservation of infinite coproducts by $G$ is assumed.

\begin{proof}
\emph{Enough projectives in the category of complexes.}
For a projective $P\in\mathcal A$, let $S^0(P)$ be $P$ in degree zero.
For $n\geq1$, let $D^n(P)$ have $P$ in degrees $n$ and $n-1$,
with the identity as its only nonzero differential.  There are natural
identifications
\[
 \operatorname{Hom}_{\mathcal B}(S^0(P),C)
       =\operatorname{Hom}_{\mathcal A}(P,C_0),
 \qquad
 \operatorname{Hom}_{\mathcal B}(D^n(P),C)
       =\operatorname{Hom}_{\mathcal A}(P,C_n).
\]
An epimorphism of complexes is an epimorphism in every degree.
Consequently these stalk and disk complexes are projective objects of
$\mathcal B$.

For any $C$, choose epimorphisms $P_n\to C_n$ with each $P_n$
projective.  The associated chain maps give an epimorphism
\[
 Q=S^0(P_0)\oplus\bigoplus_{n\geq1}D^n(P_n)\longrightarrow C.
\]
This coproduct exists without any assumption about arbitrary
coproducts in $\mathcal A$: in degree $j$ its value is the finite
biproduct $P_j\oplus P_{j+1}$.  A map from $Q$ lifts across any
epimorphism of complexes by lifting separately its maps from the
$P_n$ into degree $n$.  Hence $Q$ is projective.  Iterating on kernels
gives a projective resolution of every object of $\mathcal B$.

\emph{The finite-diagonal total-complex argument.}
We record the elementary convergence fact used below.  For a first
quadrant bicomplex $X_{r,s}$ with commuting differentials $d_h,d_v$,
give
\[
 \operatorname{Tot}(X)_N=\bigoplus_{r+s=N}X_{r,s},
 \qquad d=d_h+(-1)^r d_v.
\]
If every column is acyclic, the total complex is acyclic.  Indeed, for
a total cycle choose its component with largest first index $r$.
The cycle equation makes that component a vertical cycle, since there
is no component with first index $r+1$.  A vertical primitive, with
the appropriate sign, removes it by subtracting a total boundary;
the only newly introduced component has first index $r-1$.  Repeating
terminates after finitely many steps.  The same proof, starting with
the largest second index and using horizontal primitives, works if
every row is acyclic.  This argument uses only the finitely many
positions on a diagonal.  In a general $\mathcal A$ one may perform
this cycle calculation after applying
$\operatorname{Hom}_{\mathcal A}(P,-)$ for each projective $P$.
That functor is exact and commutes with the finite sums in question;
enough projectives makes these functors detect zero homology.

Two consequences will be used.  A morphism of first quadrant
bicomplexes that is a quasi-isomorphism in each column induces a
quasi-isomorphism of total complexes: take its columnwise mapping
cone, apply the column argument, and identify its total complex with
the mapping cone of the total map, using the usual signs.  Also, if
$X_{\bullet,s}\to C_s$ is an exact augmented horizontal complex for
each $s$, the augmentation
$\operatorname{Tot}(X)\to C$ is a quasi-isomorphism.  For this, insert
$C_s$ as the zeroth column and shift the old columns right by one;
the resulting first quadrant bicomplex has exact rows, and its total
complex is the cone of the augmentation, up to the corresponding
sign convention.  Thus its cone is acyclic.  Neither consequence
uses a totalization with infinite sums in a fixed degree.

\emph{Projective tests imply all tests.}
Assume the second condition, and write $H=\Lambda G$.  Choose a
projective resolution in the abelian category $\mathcal B$,
\[
 \cdots\longrightarrow P_2\longrightarrow P_1
       \longrightarrow P_0\longrightarrow C\longrightarrow0 .
\]
Regard $(P_r)_s$ as a first quadrant bicomplex.  Since $H$ is exact,
the augmented horizontal complex
\[
 \cdots\longrightarrow H(P_1)\longrightarrow H(P_0)
       \longrightarrow H(C)\longrightarrow0
\]
is also exact in $\mathcal B$, hence in each internal degree.
The preceding augmentation argument gives quasi-isomorphisms
$\operatorname{Tot}(P_\bullet)\to C$ and
$\operatorname{Tot}(H(P_\bullet))\to H(C)$.

Naturality of the counit gives a morphism between these bicomplexes.
On its $r$th column it is $\varepsilon_{P_r}$, a quasi-isomorphism by
assumption. The column argument therefore makes the top arrow in
the commutative square
\[
\begin{array}{ccc}
 \operatorname{Tot}(H(P_\bullet))&
       \xrightarrow{\operatorname{Tot}(\varepsilon)}&
                         \operatorname{Tot}(P_\bullet)\\
 \downarrow&&\downarrow\\
 H(C)&\xrightarrow{\varepsilon_C}&C
\end{array}
\]
a quasi-isomorphism.  Both vertical arrows are quasi-isomorphisms.
On homology the square then shows that $\varepsilon_C$ is an
isomorphism in every degree.  This proves the first condition.  The
converse implication to the second is immediate.

Notice that $H$ was applied separately to the resolution objects
$P_r$.  We neither equated $H(\operatorname{Tot}P_\bullet)$ with
$\operatorname{Tot}H(P_\bullet)$ nor applied $G$ to any infinite
coproduct.

\emph{All counits imply the localization equivalence.}
Suppose the first condition holds.  For a morphism $v:C\to C'$,
naturality gives
\[
 \varepsilon_{C'}\circ\Lambda G(v)=v\circ\varepsilon_C.
\]
Since both counits are quasi-isomorphisms, the two-out-of-three
property gives
\[
 \Lambda G(v)\in\mathcal Q
       \quad\Longleftrightarrow\quad v\in\mathcal Q.
\]
This is $G^{-1}(W)=\mathcal Q$ and, in particular, allows $G$ to
induce $\overline G$ on the localizations.  The triangular identity
\[
 \varepsilon_{\Lambda X}\circ\Lambda(\eta_X)
       =1_{\Lambda X}
\]
gives $\Lambda(\eta_X)\in\mathcal Q$, hence $\eta_X\in W$.
The unit and counit therefore descend to natural isomorphisms on the
localized categories.  Their triangular identities remain valid,
so $\overline\Lambda$ and $\overline G$ are inverse equivalences.

\emph{The converse with the localized adjunction made explicit.}
Assume the third condition.  Because both functors preserve the
respective classes being inverted, their original unit and counit
descend to natural transformations on the localizations: naturality
on original arrows implies naturality on their formal inverses.
Their triangular identities give an adjunction
$\overline\Lambda\dashv\overline G$.

The unit of an adjunction whose left adjoint is fully faithful is an
isomorphism: for each $Z$, its map on $\operatorname{Hom}(Z,-)$ is,
under adjunction, the fully faithful map on morphisms induced by the
left adjoint, so the conclusion follows from Yoneda.  Apply this to
$\overline\Lambda$.  The triangular identity makes the counit an
isomorphism on objects in its essential image, and essential
surjectivity and naturality make it an isomorphism on every object.
Thus the image of each $\varepsilon_C$ in
$\mathcal B[\mathcal Q^{-1}]$ is invertible.

Finally, a chain map whose image is invertible in this localization
is a quasi-isomorphism.  Each homology functor factors through the
localization, so it sends that map to an isomorphism.  This proves
the first condition and completes all implications.
\end{proof}

\paragraph{Scope of the projective test and of the source reading.}
The phrase \emph{projective object of
$\operatorname{Ch}_{\geq0}(\mathcal A)$} means projectivity in the
abelian category.  It is stronger than being projective in every
degree: a nonzero projective object placed only in degree one is
termwise projective but is not a projective complex.  For example the
evident epimorphism $D^1(P)\to S^1(P)$ admits no chain section when
$P\ne0$.  Testing the counit on \emph{all} termwise projective
complexes would be a stronger assumption than condition two and
therefore also sufficient; the two classes must nevertheless not be
identified.  Testing only projectives of $\mathcal A$ placed in
degree zero is not what has been proved here.

The source's arrow from $\mathcal M$ to the category of complexes is
the left adjoint; its reverse arrow is the right adjoint.
Handwritten page~53 explicitly discusses the right adjoint preserving
epimorphisms, which, together with its automatic left exactness,
gives the exactness used here.  Page~52 displays the all-counits
quasi-isomorphism condition and an equivalence after localization.
F-SM-17 establishes this precise modern formulation completely.
It does not assert that every overwritten formulation of
$(ii')$ or $(ii'')$ has been deciphered, nor does it infer a localization
equivalence from the left functor alone without the stated compatibility
of its right adjoint.

The primary locator is Grothendieck, A. (n.d.),
\emph{Pursuing stacks: The Modelizing Story, Chapter VII,
Linearization of homotopy types}, cote~134-8, handwritten
pp.~52--53, physical PDF pp.~53--54:
\url{https://grothendieck.umontpellier.fr/134-8.pdf}.
The argument above is a new complete presentation of the stated
criterion, with no claim that its categorical mechanism is historically
new or that the overwritten variants have been recovered verbatim.
